\subsection{Two scalar inputs suffice for the sharp critical law}
\label{sec:rank-two-critical}

The scalar bottleneck argument extends to a first layer of rank two.
The extra coordinate requires care. A direct tensor Bernstein
coefficient averages over a two-dimensional hypergeometric grid.
Computing that average at every selected level would have infinite
expected work. The construction below changes one coordinate operator
at a time, so every correction contains one scalar average.

\begin{theorem}[Critical sampling after a rank-two first layer]
\label{thm:rank-two-critical}
Consider a zero-bias width-$n$, depth-$D\ge1$ tanh network with
absolute row sums at most one and a first weight matrix of rank at
most two. Its weights are dyadic with at most $b$ bits. Put
$B=16+b+\lceil\log_2(nD+2)\rceil$.
There is a weight-only data structure that samples the exact output
using $O(\sqrt{D+1})$ expected input probes and
$O(\sqrt{D+1}\,B^4)$ expected bit operations. At
$D=\lceil\log_2n\rceil$ and $b=O(\log n)$, preprocessing and
storage are $O(Dn^2\log^6n)$.
For arbitrary $D,b$, preprocessing is bounded by
\[
 O\!\left(n^2B^3+D^2[n^2K^2+nK^4+nM]M^2\right),
 \qquad K=O(B),\quad M=O(B+D).
\]
The constants are absolute. The worst-case query order is
$\Theta(\sqrt D)$ when $n\ge D+1$ and
$b\ge\lceil\log_2(D+1)\rceil$.
\end{theorem}

We first give the two-coordinate factory, with the norms of
Section~\ref{sec:fixed-rank-bernstein}.

\begin{lemma}[A coordinatewise Bernstein correction]
\label{lem:rank-two-bernstein}
Let $f\in C^6([0,1]^2)$ satisfy $\|f\|_\infty\le4/5$.
For $2\le k\le6$, let nonnegative rational $M_k$ bound
$\|D^kf\|_{1,\mathrm{ent}}$. Suppose certified intervals for
$f,f_{11},f_{22},f_{1122}$ are available at rational arguments.
Put
\[
 S=M_2+M_3+M_4,\qquad U=M_4+M_5+M_6,
\]
and choose the least power of two
\[
 m_0\ge\max\{2,2M_2,\sqrt{32S},(4U)^{1/3}\}.
\]
There is an exact sign factory of mean $f(p_1,p_2)$, from two
independent Bernoulli sources, using at most $2m_0$ source bits
in expectation. Every known correction coin requires one scalar
hypergeometric average with $O(m)$ summands at level $m$.
\end{lemma}

\paragraph{Proof idea.}
Specialize the coordinatewise product telescope to two source biases. The
unchanged coordinate contributes only a second derivative, so each
correction needs four known derivatives and one scalar average.

\begin{proof}
Apply Lemma~\ref{lem:fixed-rank-bernstein} with $r=2$. Its
starting conditions are $m_0^2\ge32S$ and $m_0^3\ge4U$,
which are precisely the displayed conditions. The common correction
envelope is $\overline A=S+U/(8m_0)\le m_0^2/16$.
The product $(I-A_1/m)(I-A_2/m)$ uses only
$f,f_{11},f_{22},f_{1122}$. The same is true of each coordinate
correction, with its scalar hypergeometric average. The general
source bound becomes $2m_0$. Integer-power comparisons select the
starting level without testing equality of real roots.
\end{proof}

The finite-bit implementation uses a complex strip. Suppose the
first preactivations are $a_iz$, with $|a_i|\le1$, and every later
row has absolute sum at most one. Write $F(z)$ for the final
activation. A rank-one first layer has this form: choose a row of
largest absolute sum as the scalar source row. The strip estimate
also applies to several scalar inputs after bounding the imaginary
part of their first linear combination.

\begin{lemma}[A uniform complex strip]
\label{lem:signedstrip}\label{sec:signed-bottleneck}
Every scalar neuron, and in particular $F$, is holomorphic on
$|\Im z|<1/(4\sqrt D)$. Throughout this strip every neuron output
has modulus at most two.
\end{lemma}
\paragraph{Proof idea.}
The imaginary part of tanh is bounded by a scalar tangent recurrence. Its
reciprocal-square bound keeps every preactivation away from the poles
throughout the depth.

\begin{proof}
For real $u,v$ with $|v|<\pi/4$,
\[
 |\Im\tanh(u+iv)|
 \le\frac{|\sin(2v)|}{1+\cos(2v)}=\tan|v|,
 \qquad |\Re\tanh(u+iv)|\le1.
\]
The denominator is positive, so there is no pole.
Also $\tan t\le t/\sqrt{1-t^2}$ for $0\le t<1$, by
$\sin^2t\le t^2$. Absolute row sums at most one show inductively
that, at input imaginary magnitude $y$, every layer-$\ell$
imaginary magnitude is at most
\[
 \frac{|y|}{\sqrt{1-\ell y^2}}.
\]
For $|y|\le1/(4\sqrt D)$ this is at most $1/\sqrt{15D}<1/2$.
Consequently all intermediate complex preactivations avoid the
poles of tanh, establishing holomorphy by composition. Their real
outputs have absolute value at most one and their imaginary outputs
have absolute value below $1/2$, proving the modulus bound two.
The argument uses absolute weights and does not assume positivity.
\end{proof}

\paragraph{Proof idea for Theorem~\ref{thm:rank-two-critical}.}
Two basis rows turn the first layer into two unknown scalar means. The
composite network is a known function on their enlarged square. Its output
radius decays as $D^{-1/2}$, while its normalized second, fourth, and sixth
derivative bounds grow as $D,D^2,D^3$. The lemma therefore costs $O(D)$
conditionally on a gate of probability $O(D^{-1/2})$. A two-dimensional
Taylor table supplies every known coefficient request within the
preprocessing budget.

\begin{proof}
A zero first matrix gives a fair output. A single layer is handled
by the existing unit-row factory. We henceforth assume $D\ge2$ and
first-layer rank $r\in\{1,2\}$.
Find $r$ independent rows and $r$ columns on which their minor is
nonsingular. If a row's representation in the current row basis
has a coefficient of absolute value greater than two, replace that
basis row by this row. The fixed-column determinant grows by a
factor greater than two. Every nonzero such determinant is at least
$2^{-br}$ and at most $r!$ in absolute value. Thus at most
$br+\lceil\log_2(r!)\rceil+1$ replacements occur. Exact elimination
and row checks cost $O(n^2B^3)$ bits for $r\le2$.

The resulting basis rows $v_a$ have absolute sums at most one, and
every first-layer row is $\sum_a a_{ia}v_a$ with $|a_{ia}|\le2$.
Their coefficients are ratios of fixed-order dyadic minors and have
$O(b)$ bits. Fresh weighted-coordinate draws supply independent
signs of means $z_a=v_a\cdot x$, each with at most one input probe.
The network output is now a known function $F(z_1,z_2)$; a missing
second coordinate is ignored. On the enlarged square, first-layer
preactivations have magnitude at most $c=4$. First-layer outputs have
magnitude at most one, and subsequent rows remain contractive.
Consequently
\[
 |F|\le\tanh^{\circ(D-1)}(1)
       \le\rho_D:=\sqrt{3/(2D+1)}.
\]
If $2\rho_D\ge1$, set $q=1$. Otherwise let $q$ be the smallest
number of the form $2^{-j}$, $j\ge0$, with $q\ge2\rho_D$.
Comparisons of squares determine this dyadic number exactly.
Then
\[
 q\le6/\sqrt{D+1},\qquad q^{-1}\le\sqrt D,
 \qquad \|F/q\|_\infty\le4/5.
\]
The range bound for $q=1$ uses $D\ge2$ and $\tanh1<4/5$.

Here are explicit derivative constants. Put $T_k=2^kk!$ and
$C_1=1$, and define
\[
 C_k=T_k+4\sum_{\nu=2}^kT_\nu
       \mathsf B_{k,\nu}(C_1,\ldots,C_{k-\nu+1}),
 \qquad k\ge2,
\]
where $\mathsf B_{k,\nu}$ is the partial exponential Bell
polynomial, with coefficient
$k!/\prod_j(m_j!(j!)^{m_j})$ for
$\sum_jm_j=\nu$ and $\sum_jjm_j=k$.
The derivative polynomial of tanh has coefficient sum at most $T_k$.
For even $k$ it is odd, giving the bound
$T_k|\tanh u|$; for odd $k$ the bound is $T_k$.
Use $|\tanh'|\le1$ for the first derivative.

Let $E_{k,\ell}$ be the largest uniform total derivative norm at
layer $\ell$ on the enlarged square. The first layer has
$E_{k,1}\le T_kc^k$. Later scalar chain rules and row contraction
give
\[
 E_{k,D}\le C_kc^kD^{(k-1)/2}.
\]
To check this induction, the radius at layer $\ell$ is at most
$2/\sqrt\ell$. A nonlinear Bell term with $\nu\ge2$ factors
is bounded by
$2T_\nu\mathsf B_{k,\nu}(C)c^k\ell^{(k-3)/2}$.
Even $\nu$ uses the radius factor; odd $\nu\ge3$ supplies the
same power from its number of factors. Finally,
$\sum_{\ell=2}^D\ell^{(k-3)/2}\le2D^{(k-1)/2}$.
This proves the displayed recurrence, including its first-layer term.
The induction uses only the bound $c$ on the first preactivations and
holds at every order $k$.

For $f(p)=F(2p-1)/q$, we may therefore take
\[
 M_k=2^kC_kc^k(D+1)^{\lceil k/2\rceil},\qquad 2\le k\le6.
\]
In Lemma~\ref{lem:rank-two-bernstein}, these give
$m_0=O(D+1)$ and $\overline A=O((D+1)^2)$. Open the gate $q$
and run that factory on its open branch, returning a fair sign
otherwise. Its expected input count is $O(\sqrt{D+1})$.

We provide the finite-bit implementation. The complex-strip argument
of Lemma~\ref{lem:signedstrip} applies when
$\max_a|\Im z_a|<1/(16\sqrt D)$, since the first affine row has
absolute sum at most four. Choose dyadic
$1/(256\sqrt D)<a\le1/(128\sqrt D)$ and grid centers on the
square at spacing $a$. There are $O(D)$ centers. In scaled variables
$z_a=c_a+2au_a$, the polydisk $|u_a|\le2$ lies strictly inside
the tube. Every neuron has modulus at most two, so its coefficient
of $u^\nu$ is bounded by $2^{1-|\nu|}$.

Let $H=n^2D$, choose a dyadic $\varepsilon=2^{-P}\le H^{-2}$
with $P=O(B)$, and put
\[
 S_0=P+3\lceil\log_2(D+1)\rceil+64,\qquad
 K=16(P+\lceil\log_2(D+1)\rceil+64).
\]
Store the total-degree-$K$ expansions with coefficient uncertainty
at most $2^{-S_0}$. On the half-radius polydisk, coefficient errors
in the mixed fourth derivative sum to at most $256\,2^{-S_0}$.
Its scale conversion is at most
$1/(qa^4)\le2^{32}D^{5/2}$. The geometric tail before conversion
is at most $2^{12}(K+2)^5 4^{-K}$. These choices provide certified
width-$\varepsilon$ intervals for $f,f_{11},f_{22},f_{1122}$.
There are $O(K^2)$ coefficients. Rounded Horner evaluation, with
$O(\log(K+D+2))$ guard bits, costs
$O((B+\operatorname{bits}(p))^4)$.

We now bound the finite precision of the two-variable scaled-series
computation. In the total-degree quotient
algebra, the coefficient norm is submultiplicative. Exponentials
satisfy the Euler recurrence
\[
 |\nu|e_\nu=\sum_{0<\mu\le\nu}|\mu|a_\mu e_{\nu-\mu}.
\]
Write $\mathcal E=\sum_i u_i\partial_i$. If a computed exponential
$\widetilde E$ has residual
$R=\mathcal E\widetilde E-(\mathcal EA)\widetilde E$, then
$\mathcal E(e^{-A}\widetilde E)=e^{-A}R$ in the truncated algebra.
The inverse of $\mathcal E$ on series with zero constant term
divides a coefficient by its total degree and has norm at most one.
Multiplication by $e^{-A}$, this integration, and multiplication by
$e^A$ therefore bound the error, including the separately enclosed
constant coefficient. There is no additional degree factor in this
inverse operation. Each complex preactivation has imaginary part
below one half, so $\Re e^{2s}>0$ and $|1+e^{2s}|\ge1$.
Cauchy's bound consequently gives coefficient norms eight
for neurons, $4e^{10}$ for exponentials, and four for reciprocals.
The first preactivation is affine with coefficient norm below five;
later preactivations have norm at most eight. A perturbed exponential
input therefore has norm at most seventeen. The exponential and
inverse identities bound layer amplification by $2^{40}$ and local
error by $2^{90}n(K+2)^8 2^{-M}$. One sufficient choice is
\[
 M=S_0+128D+\lceil\log_2(256nD(K+2)^{12})\rceil+4b+256.
\]
It verifies all perturbation margins and leaves error below
$2^{-S_0-16}$. Bounded rounded intermediates follow by applying the
same residual identity in every further quotient that discards
terms above a specified total degree. Induction over that degree
keeps the computed coefficient norm within one of the exact bound;
all numerator partial sums are then polynomially bounded in $K$.
This closes the intermediate-size assumption in the rounding count.
Dense rows cost $O(n^2K^2)$ coefficient operations,
nonlinear recurrences $O(nK^4)$, and constant exponentials
$O(nM^3)$ bit operations per layer and center. This proves the
preprocessing bound in the theorem and the stated storage budget.

Each selected correction requires one mode-normalized
hypergeometric average of $O(m)$ oracle values. The recurrence in
Theorem~\ref{thm:implicitrankone} evaluates it with
$O(\log(m+2))$ guard bits. Its cheap cost is
$O(m(B+\log(m+2))^4)$, and its absolute uncertainty before dividing
by its branch mass is $O(\varepsilon)$. The base uses the same
oracle without an average.

An augmented forward pass for derivatives through any fixed order
supplies arbitrary further precision in
$O(H(B+\operatorname{bits}(p)+t)^4)$ bit operations, with a constant
depending on that order. Here is an explicit guard-bit argument.
If each completed derivative update has rounding error at most $\zeta$,
write $e_{k,j}$ for the largest error of order $k$ at layer $j$.
The coefficient of the highest derivative is $\tanh'$ at the
approximate real preactivation and has magnitude at most one.
Every remaining chain-rule term is a polynomial in lower derivatives.
The derivative bounds already proved, with a one-unit perturbation
margin, therefore give for every fixed $k$
\[
 e_{0,j+1}\le e_{0,j}+\zeta,\qquad
 e_{k,j+1}\le e_{k,j}
   +A_k(D+1)^k\sum_{a<k}e_{a,j}+\zeta.
\]
Induction on $k$ and summation over $j\le D$ give
$e_{k,j}\le A'_k\zeta(D+1)^{(k+1)^2}$.
Choosing $\zeta$ below the target accuracy divided by this polynomial
also verifies the perturbation margin. Scale conversion by $q$ and
the dyadic cell radius costs only further powers of $D$.
Assigning each row operation its share of $\zeta$ adds $O(\log n)$
guard bits. Thus all guard lengths are logarithmic in $D,n$.
On a branch of mass $c_m$, the unresolved uniform interval has measure
$O(\min\{1,\varepsilon/c_m\})$. For $c_m=O(\varepsilon)$,
start direct evaluation at absolute precision proportional to $c_m$.
For a cached miss, start at absolute precision proportional to
$\varepsilon$. After $t$ extra accuracy bits, the unresolved set of
uniform values has measure $O((\varepsilon/c_m)2^{-t})$.
Multiplying by the branch mass before summing refinement stages
therefore charges $O(Hm\min\{c_m,\varepsilon\}
(B+\log(m+2)+P)^4)$ to that branch. This uses unconditional interval
measures and imposes no independence between work and the miss event.
The total expensive refinement work is consequently bounded by
\[
 O\!\left(
 H\sum_{m=2^\ell m_0}
    \min\{\overline A/m^2,\varepsilon\}
       m(B+\log(m+2)+P)^4\right)
 =O\!\left(H\sqrt{\overline A\varepsilon}\,B^4\right)
 =O((D+1)B^4).
\]
The cheap averages cost
$O((\overline A/m_0)B^4)=O((D+1)B^4)$, including their
geometric logarithmic moments. Source choices and counters fit this
bound. Multiplication by the gate proves the online bit estimate.
All comparisons terminate almost surely, all selected levels are
finite, and the exact telescope proves the target law.

At logarithmic depth and parameter precision, the dense preprocessing
term is $O(n^2\log^6n)$ and the nonlinear term is
$O(n\log^8n)$; both fit $O(Dn^2\log^6n)$. The rank-one critical
mean-chain lower bound is contained in this class, proving the final
claim under the stated width and grid conditions.
\end{proof}

Rank one is a special case of rank two with a simpler basis.

\begin{corollary}[A sharp critical law with a rank-one first layer]
\label{thm:signed-bottleneck}
For a zero-bias width-$n$, depth-$D$ tanh network with row norms
at most one and first-layer rank at most one, exact sampling uses
$O(\sqrt{D+1})$ expected input probes and
$O(\sqrt{D+1}\,B^4)$ expected online bit operations, with $B$ as in
Theorem~\ref{thm:rank-two-critical}. At $D=\lceil\log_2n\rceil$
and $b=O(\log n)$, preprocessing and storage are $O(Dn^2\log^6n)$.
Hidden weights may have either sign. The query order is sharp when
$n\ge D+1$ and $b\ge\lceil\log_2(D+1)\rceil$.
\end{corollary}
\paragraph{Proof idea.}
One scalar source describes the first layer, and a mean chain supplies
the matching lower bound.

\begin{proof}
Apply Theorem~\ref{thm:rank-two-critical} with first-layer rank one.
Its basis construction reduces here to choosing a nonzero row of
largest absolute sum. The mean-chain lower bound in
Theorem~\ref{thm:finitecritical} has a rank-one first matrix and
fits the stated grid.
\end{proof}
